\documentclass[11pt,letterpaper]{article}

% ── packages ──────────────────────────────────────────────────────────────────
\usepackage[margin=1in]{geometry}
\usepackage{amsmath,amssymb}
\usepackage{graphicx}
\usepackage{booktabs}
\usepackage{multirow}
\usepackage{caption}
\usepackage{subcaption}
\usepackage{hyperref}
\usepackage{xcolor}
\usepackage{enumitem}
\usepackage{float}
\usepackage{microtype}
\usepackage[numbers,sort&compress]{natbib}

\hypersetup{colorlinks=true, linkcolor=blue!70!black,
            citecolor=blue!70!black, urlcolor=blue!70!black}
\captionsetup{font=small, labelfont=bf}
\setlength{\parskip}{4pt}

% ── title ─────────────────────────────────────────────────────────────────────
\title{%
  \Large\textbf{Cancer Subtype Discovery and Classification}\\[4pt]
  \large from TCGA Pan-Cancer Gene Expression Data\\[6pt]
  \normalsize 02-620 Machine Learning for Scientists --- Spring 2026
}
\author{Tolga Ozgun \and Yuchen Shao \and Roy Hung}
\date{April 2026}

\begin{document}
\maketitle

% ─────────────────────────────────────────────────────────────────────────────
\begin{abstract}
We present an end-to-end machine learning pipeline for cancer subtype
discovery and classification applied to bulk RNA-seq data from The Cancer
Genome Atlas (TCGA) Pan-Cancer dataset.  Starting from raw HTSeq-FPKM-UQ
expression matrices across six cancer types (BRCA, LUAD, KIRC, PRAD, COAD,
GBM), we implement three methods entirely from scratch: (1)~unsupervised
clustering via PCA + K-Means, (2)~supervised one-vs-rest logistic regression
with $\ell_2$-regularisation, and (3)~a multi-layer perceptron (MLP) in
PyTorch.  All supervised models are evaluated using 5-fold stratified
cross-validation with preprocessing fit inside each fold to prevent data
leakage.  K-Means achieves ARI\,=\,0.794, NMI\,=\,0.849; logistic regression
reaches macro-F1\,=\,0.9927 and the MLP reaches macro-F1\,=\,0.9981
(5-fold CV: $0.9983 \pm 0.0010$) on the held-out test set.  Motivated by TA
feedback, we extend the pipeline with biological validation via feature
importance analysis: LR weights are projected from PCA space back to gene
space, revealing well-established tissue-specific markers (e.g., NKX2-1/TTF-1
for LUAD, PCA3 for PRAD, GATA3 for BRCA) as the top discriminative genes.
MLP permutation importance confirms that the first eight principal components
together account for the vast majority of classification signal, with a sharp
drop-off thereafter.  These findings confirm that the models learn biologically
meaningful signals rather than technical artefacts.
\end{abstract}

% ─────────────────────────────────────────────────────────────────────────────
\section{Introduction}
\label{sec:intro}

Cancer is not a single disease.  Even within the same organ, tumours can differ
dramatically in molecular profile, prognosis, and response to therapy.  Bulk
RNA-sequencing captures a genome-wide snapshot of gene expression in a tumour
sample, providing a high-dimensional view of its molecular state.  The
challenge we address is threefold:

\begin{enumerate}[leftmargin=2em,itemsep=1pt]
  \item \textbf{Unsupervised discovery.}  Can meaningful cancer subtypes be
        recovered from unlabelled gene expression data using PCA followed by
        K-Means clustering?
  \item \textbf{Linear classification.}  How accurately can a multi-class
        logistic regression model predict cancer type from PCA-reduced features?
  \item \textbf{Non-linear classification.}  Does a deep neural network offer
        a measurable gain over the linear baseline?
\end{enumerate}

A fourth question --- raised during TA review --- drives the extension in
Section~\ref{sec:importance}:

\begin{enumerate}[leftmargin=2em,itemsep=1pt]
  \setcounter{enumi}{3}
  \item \textbf{Biological validity.}  Do the models learn tissue-specific gene
        expression signatures, or could the high accuracy reflect technical
        batch structure that survived normalisation?
\end{enumerate}

% ─────────────────────────────────────────────────────────────────────────────
\section{Dataset}
\label{sec:dataset}

\subsection{Source}

We use the TCGA Pan-Cancer (PANCAN) gene expression dataset downloaded via the
GDC Data Portal API\footnote{\url{https://portal.gdc.cancer.gov/}}.
Specifically:

\begin{itemize}[leftmargin=2em,itemsep=1pt]
  \item \textbf{Expression matrix} (UUID
        \texttt{3586c0da-64d0-4b74-a449-5ff4d9136611}):
        HTSeq-FPKM-UQ values, batch-corrected with EBPlusPlus across all TCGA
        projects.
  \item \textbf{Clinical metadata} (UUID
        \texttt{1b5f413e-a8d1-4d10-92eb-7c4ae739ed81}):
        patient-level cancer type annotations from the TCGA Clinical Data
        Resource.
\end{itemize}

\subsection{Cancer Types and Class Distribution}

We focus on six cancer types spanning distinct tissue origins
(Table~\ref{tab:dataset} and Figure~\ref{fig:class_dist}).

\begin{table}[H]
  \centering
  \caption{Dataset composition after preprocessing.}
  \label{tab:dataset}
  \begin{tabular}{llr}
    \toprule
    Cancer Type & Abbreviation & Samples \\
    \midrule
    Breast Invasive Carcinoma         & BRCA & 1{,}215 \\
    Kidney Renal Clear Cell Carcinoma & KIRC &   606 \\
    Lung Adenocarcinoma               & LUAD &   576 \\
    Prostate Adenocarcinoma           & PRAD &   550 \\
    Colon Adenocarcinoma              & COAD &   492 \\
    Glioblastoma Multiforme           & GBM  &   168 \\
    \midrule
    \textbf{Total}                    &      & \textbf{3,607} \\
    \bottomrule
  \end{tabular}
\end{table}

\begin{figure}[H]
  \centering
  \includegraphics[width=0.72\textwidth]{../results/runs/run_20260419_175643/graphs/fig01_class_distribution.png}
  \caption{Sample count per cancer type.  BRCA is the largest class (1,215
           samples) and GBM the smallest (168 samples), introducing mild class
           imbalance addressed by macro-averaged metrics.}
  \label{fig:class_dist}
\end{figure}

% ─────────────────────────────────────────────────────────────────────────────
\section{Data Preprocessing}
\label{sec:preprocessing}

\paragraph{Loading and filtering.}
We first read only the header of the 1.8\,GB TSV to identify columns matching
our six cancer types, then load only those columns --- reducing I/O substantially.

\paragraph{Log-transformation and NaN handling.}
Raw FPKM-UQ values are $\log_2(x+1)$ transformed to stabilise variance.
Missing values are filled with 0.

\paragraph{Variance filtering.}
Genes with variance $<10^{-6}$ are removed, leaving \textbf{20,297 gene
features}.

\paragraph{Train / test split.}
15\% of samples are held out as a final test set via stratified random
splitting (seed 42):
\begin{itemize}[leftmargin=2em,itemsep=0pt]
  \item Development set: 3,065 samples (85\%)
  \item Test set (held out): 542 samples (15\%)
\end{itemize}

\paragraph{Standardisation and PCA.}
A single \texttt{StandardScaler} and PCA (50 components) are fit on the
development set and applied to the test set.  For cross-validation, both
transformers are re-fit inside each fold's training partition only, so no
information from the validation partition influences the fitted transformers
(preventing data leakage).

% ─────────────────────────────────────────────────────────────────────────────
\section{Methodology and Implementation}
\label{sec:methodology}

All three ML methods are implemented from scratch in NumPy (Methods 1 \& 2)
or PyTorch (Method 3).  Scikit-learn is used only for evaluation metrics,
data splitting, and preprocessing utilities.

\subsection{Method 1: PCA + K-Means Clustering}

\paragraph{PCA.}
We implement PCA via truncated SVD: given centered data
$\mathbf{X} \in \mathbb{R}^{n \times p}$, we compute
$\mathbf{X} = \mathbf{U}\mathbf{S}\mathbf{V}^\top$ and retain the top $k=50$
right singular vectors.  Explained variance ratios are computed as
$r_i = \lambda_i / \sum_j \lambda_j$ where $\lambda_i = s_i^2 / (n-1)$, with
the denominator summed over \emph{all} components so the ratios are properly
normalised.  The top 50 PCs capture \textbf{63.7\%} of total variance.

\paragraph{K-Means.}
Lloyd's algorithm with $k{=}6$, $n\_\text{init}{=}10$ random restarts, and
convergence tolerance $10^{-4}$.  After convergence, labels are re-assigned
using the final centroids before computing inertia (correcting a subtle bug
where labels could refer to pre-update centroids).

\subsection{Method 2: Logistic Regression (One-vs-Rest)}

We train $C=6$ binary logistic classifiers with $\ell_2$-regularised
cross-entropy loss:

\begin{equation}
  \mathcal{L}(\mathbf{w}, b) =
    -\frac{1}{n}\sum_{i=1}^n \Bigl[\tilde{y}_i \log \hat{y}_i +
      (1-\tilde{y}_i)\log(1-\hat{y}_i)\Bigr]
    + \frac{\lambda}{2n} \|\mathbf{w}\|_2^2
\end{equation}

Training uses gradient descent with $\eta{=}0.1$ for 1,000 iterations.
Hyperparameter search over $\lambda \in \{0.001, 0.01, 0.1\}$ is performed via
5-fold stratified CV, with preprocessing fit inside each fold.

\subsection{Method 3: Multi-Layer Perceptron (MLP)}

Feedforward network implemented in PyTorch:
\begin{equation*}
  \underbrace{50}_{\text{input}} \to
  \underbrace{512}_{\text{Linear+ReLU+BN+Dropout}} \to
  \underbrace{256}_{\text{Linear+ReLU+BN+Dropout}} \to
  \underbrace{6}_{\text{output}}
\end{equation*}
Trained with Adam ($\eta{=}0.001$), batch size 64, 50 epochs,
\texttt{CrossEntropyLoss}.  Evaluated with the same 5-fold CV framework as
the LR model.  To ensure reproducibility and prevent any form of test-set
influence during final training, a fixed random seed (\texttt{torch.manual\_seed(42)})
is set before each training run, and a 10\% monitor split is carved from the
development set for early-stopping diagnostics (the test set is never used
as a validation signal).

\subsection{Evaluation Metrics}

\textbf{Clustering:} ARI, NMI, Silhouette score.
\textbf{Classification:} Macro-averaged F1, Precision, Recall; per-class
confusion matrix; 5-fold CV mean $\pm$ std F1.

% ─────────────────────────────────────────────────────────────────────────────
\section{Results}
\label{sec:results}

\subsection{PCA Analysis}

The top 50 PCs explain 63.7\% of total variance
(Figure~\ref{fig:pca_variance}).  The scree plot shows a characteristic
elbow around PC~8--10; subsequent components each contribute $<0.5\%$.
Even in two dimensions (Figure~\ref{fig:pca_labels}), the six cancer types
form largely non-overlapping clusters --- a strong indicator that supervised
classification will be straightforward.

\begin{figure}[H]
  \centering
  \includegraphics[width=\textwidth]{../results/runs/run_20260419_175643/graphs/fig02_pca_variance.png}
  \caption{PCA explained variance on the development set.
           \textit{Left:} per-component variance (scree plot).
           \textit{Right:} cumulative variance; the dashed line marks the
           63.7\% explained by the top 50 components.}
  \label{fig:pca_variance}
\end{figure}

\begin{figure}[H]
  \centering
  \includegraphics[width=0.72\textwidth]{../results/runs/run_20260419_175643/graphs/fig03_pca_true_labels.png}
  \caption{PCA projection (PC1 vs PC2) coloured by true cancer type.  Even
           in two dimensions, the six cancer types are largely separable.}
  \label{fig:pca_labels}
\end{figure}

\subsection{Method 1: Unsupervised Clustering}

\begin{table}[H]
  \centering
  \caption{K-Means clustering evaluation (development set).}
  \label{tab:clustering}
  \begin{tabular}{lr}
    \toprule
    Metric & Score \\
    \midrule
    Adjusted Rand Index (ARI)        & 0.794 \\
    Normalised Mutual Information (NMI) & 0.849 \\
    Silhouette Score                 & 0.228 \\
    \bottomrule
  \end{tabular}
\end{table}

ARI\,=\,0.794 indicates strong unsupervised recovery of the true partitioning.
The Silhouette score (0.228) is consistent with published results on
high-dimensional transcriptomic data: TCGA bulk RNA-seq studies routinely
report Silhouette values in the 0.15--0.35 range even for well-validated
subtypes~\cite{weinstein2013}, because within-cluster geometric variance in
50-D PCA space is inherently large.  ARI and NMI are more informative here as
they measure agreement with known labels directly.

Figure~\ref{fig:kmeans} shows that five of the six clusters align almost
exclusively with a single cancer type.  GBM (168 samples, the smallest class)
is the primary source of residual confusion: at this sample size it cannot
reliably form a well-separated centroid in PCA space, and its cluster partially
overlaps neighbouring tissue types.  The cluster purity heatmap
(Figure~\ref{fig:purity}) confirms that BRCA, COAD, KIRC, LUAD, and PRAD each
map cleanly to dedicated clusters.

\begin{figure}[H]
  \centering
  \includegraphics[width=\textwidth]{../results/runs/run_20260419_175643/graphs/fig04_kmeans_vs_truth.png}
  \caption{\textit{Left:} samples coloured by true cancer type.
           \textit{Right:} K-Means cluster assignments; crosses mark centroids.
           The panels are visually similar, confirming ARI\,=\,0.794.}
  \label{fig:kmeans}
\end{figure}

\begin{figure}[H]
  \centering
  \includegraphics[width=0.72\textwidth]{../results/runs/run_20260419_175643/graphs/fig05_cluster_purity.png}
  \caption{Cluster purity heatmap.  Colours encode the proportion of each
           cluster belonging to each cancer type; cell numbers show raw
           counts.}
  \label{fig:purity}
\end{figure}

\subsection{Methods 2 \& 3: Supervised Classification}

\paragraph{Cross-validation.}
Figure~\ref{fig:cv_folds} shows per-fold macro-F1.  LR achieves
$0.9968 \pm 0.0027$ and MLP achieves $0.9983 \pm 0.0010$ --- both
remarkably stable across folds, with MLP showing lower variance.

\begin{figure}[H]
  \centering
  \includegraphics[width=0.75\textwidth]{../results/runs/run_20260419_175643/graphs/fig06_cv_folds.png}
  \caption{5-fold CV macro-F1.  Dashed lines show cross-fold means.}
  \label{fig:cv_folds}
\end{figure}

\paragraph{Test set evaluation.}
Table~\ref{tab:classification} summarises test-set performance.
Confusion matrices (Figures~\ref{fig:lr_cm} and~\ref{fig:mlp_eval}) and
per-class F1 (Figure~\ref{fig:perclass}) show near-perfect classification
for all cancer types.  Under LR, GBM is the only class with a slightly lower
F1 (0.96); the MLP resolves most of these boundary cases, achieving near-perfect
classification across all six types.

\begin{table}[H]
  \centering
  \caption{Test set results (542 samples, 15\% hold-out).
           LR hyperparameter: $\lambda{=}0.001$.}
  \label{tab:classification}
  \begin{tabular}{lcccc}
    \toprule
    Model & CV F1 (mean $\pm$ std) & Test F1 & Test Precision & Test Recall \\
    \midrule
    Logistic Regression & $0.9968 \pm 0.0027$ & 0.9927 & 0.9877 & 0.9982 \\
    MLP [512, 256]      & $0.9983 \pm 0.0010$ & 0.9981 & 0.9981 & 0.9982 \\
    \bottomrule
  \end{tabular}
\end{table}

\begin{figure}[H]
  \centering
  \includegraphics[width=0.62\textwidth]{../results/runs/run_20260419_175643/graphs/fig07_lr_confusion.png}
  \caption{Logistic Regression confusion matrix on the held-out test set.}
  \label{fig:lr_cm}
\end{figure}

\begin{figure}[H]
  \centering
  \includegraphics[width=0.72\textwidth]{../results/runs/run_20260419_175643/graphs/fig08_mlp_learning_curves.png}
  \caption{MLP training dynamics.  Loss decreases smoothly; validation
           accuracy converges to $\approx$99.8\% within $\sim$20 epochs.}
  \label{fig:mlp_curves}
\end{figure}

\begin{figure}[H]
  \centering
  \includegraphics[width=0.62\textwidth]{../results/runs/run_20260419_175643/graphs/fig09_mlp_confusion.png}
  \caption{MLP confusion matrix on the held-out test set.}
  \label{fig:mlp_eval}
\end{figure}

\begin{figure}[H]
  \centering
  \includegraphics[width=0.78\textwidth]{../results/runs/run_20260419_175643/graphs/fig10_perclass_f1.png}
  \caption{Per-class F1 on the test set.  GBM is the only class where LR
           makes errors; the MLP substantially reduces this gap, achieving
           near-perfect classification across all classes.}
  \label{fig:perclass}
\end{figure}

The MLP achieves test macro-F1\,=\,0.9981, a 0.54 percentage point gain over
LR (0.9927).  This advantage is also visible in CV: the MLP's cross-validation
standard deviation is nearly three times smaller ($0.0010$ vs $0.0027$),
indicating more consistent generalisation.  The consistent advantage over LR
across runs suggests a mild non-linear benefit at class boundaries, most likely
at the BRCA--PRAD interface visible in Figure~\ref{fig:pca_labels}.

% ─────────────────────────────────────────────────────────────────────────────
\section{Biological Validation via Feature Importance}
\label{sec:importance}

\emph{This section was added in response to TA feedback, which noted that
high classification accuracy alone does not rule out the possibility that
models are exploiting residual technical batch structure.  Feature importance
analysis provides a mechanism-level check: if the top discriminative genes are
known tissue-specific markers, the models are learning biology.}

\subsection{LR: Projecting Weights to Gene Space}

Each OvR logistic regression classifier produces a weight vector
$\mathbf{w}_c \in \mathbb{R}^{50}$ in PCA space.  We project these weights
back to gene space via:
\begin{equation}
  \tilde{\mathbf{w}}_c = \mathbf{V}_k^\top \mathbf{w}_c
  \quad \in \mathbb{R}^{20{,}297}
\end{equation}
where $\mathbf{V}_k \in \mathbb{R}^{20{,}297 \times 50}$ is the matrix of
PCA loading vectors.  The resulting $\tilde{\mathbf{w}}_c$ assigns a
discriminative score to every gene for class $c$, with sign indicating whether
the gene is positively (higher in this cancer type) or negatively associated.

Figure~\ref{fig:lr_genes} shows the top 15 genes per cancer type.
Table~\ref{tab:top_genes} lists the top 5 for each class and their biological
interpretation.

\begin{table}[H]
  \centering
  \caption{Top 5 discriminative genes per cancer type from LR weight
           projection, with known biological roles.}
  \label{tab:top_genes}
  \small
  \begin{tabular}{llp{7cm}}
    \toprule
    Cancer & Top genes & Biological relevance \\
    \midrule
    BRCA & THRSP, SCGB2A2, ADIPOQ, LMX1B, GATA3
         & GATA3 is a master transcription factor for luminal breast cancer
           identity~\cite{cimino2013}; SCGB2A2 (Mammaglobin) is a
           clinically used breast tissue-specific biomarker. \\[3pt]
    COAD & LY6G6D, C13orf18, TDGF3, FAM55A, NAT2
         & NAT2 (N-acetyltransferase 2) is highly expressed in colon
           epithelium and is a known colorectal cancer susceptibility gene. \\[3pt]
    GBM  & C21orf131, C2orf80, LOC643387, LOC285370, BTBD17
         & Top genes include uncharacterised loci and brain-enriched transcripts.
           GBM's smaller sample size (168 samples) reduces reliability of
           individual gene estimates relative to the larger cohorts. \\[3pt]
    KIRC & SLC22A2, AGXT2, SLC17A1, SLC5A10, SLC17A3
         & SLC22A2 (OCT2) is expressed in kidney proximal tubule cells, the
           cell of origin for clear cell RCC; the other top genes are also
           solute carrier transporters with kidney-restricted expression. \\[3pt]
    LUAD & SFTA3, SFTPA1, NKX2-1, SCGB3A2, SFTPB
         & NKX2-1 (TTF-1) is the canonical lung adenocarcinoma marker used
           clinically to confirm primary lung origin; SFTPA1 and SFTPB encode
           pulmonary surfactant proteins produced by type II pneumocytes. \\[3pt]
    PRAD & SLC45A3, PCA3, POTEG, BEND4, LMAN1L
         & PCA3 is a prostate-specific non-coding RNA that is FDA-cleared
           as a urine biomarker for prostate cancer detection; SLC45A3 is a
           prostate-restricted androgen-regulated transporter. \\
    \bottomrule
  \end{tabular}
\end{table}

\begin{figure}[H]
  \centering
  \includegraphics[width=\textwidth]{../results/runs/run_20260419_175643/graphs/fig11_lr_gene_importance.png}
  \caption{Top 15 discriminative genes per cancer type from LR weights
           projected to gene space.  Red bars indicate genes positively
           associated with each cancer type; blue bars indicate genes
           negatively associated.}
  \label{fig:lr_genes}
\end{figure}

The biological coherence of the recovered genes is strong: five of the six
cancer types yield multiple well-established clinical markers in their top 5
(NKX2-1/TTF-1 for LUAD, PCA3 for PRAD, GATA3 and Mammaglobin for BRCA, OCT2
and related transporters for KIRC).  This confirms that the classifiers are
responding to genuine tissue-of-origin expression signatures rather than
technical batch structure.

\subsection{PC Loadings: What Do the Principal Components Capture?}

Figure~\ref{fig:pc_loadings} shows the top-loading genes for PC1 and PC2 ---
the two axes that already visually separate the cancer types in the PCA scatter
plot.

\begin{figure}[H]
  \centering
  \includegraphics[width=\textwidth]{../results/runs/run_20260419_175643/graphs/fig12_pc_loadings.png}
  \caption{Top 15 genes by absolute loading magnitude for PC1 (left) and PC2
           (right).  Red bars indicate positive loadings; blue bars negative.}
  \label{fig:pc_loadings}
\end{figure}

The PC loading analysis complements the LR gene importance: while LR weights
reflect the classifier's discriminative strategy for each class, PC loadings
reveal what biological axes of variation the data itself is organised along,
independent of any labels.

\subsection{MLP Permutation Importance}

For each of the 50 PCA dimensions, we permute that column in the test set
(averaged over 3 repeats) and record the drop in macro-F1 relative to the
unpermuted baseline.  A large drop indicates the model relies heavily on that
dimension; near-zero drop indicates the dimension carries little discriminative
signal for the MLP.

\begin{figure}[H]
  \centering
  \includegraphics[width=\textwidth]{../results/runs/run_20260419_175643/graphs/fig13_mlp_permutation_importance.png}
  \caption{MLP permutation importance: macro-F1 drop when each PC dimension
           is permuted (sorted by importance; top 10 highlighted in red).}
  \label{fig:perm_importance}
\end{figure}

The results (Table~\ref{tab:perm}) reveal a clear hierarchy: PC2 and PC1 are
by far the most critical ($\Delta F1 = 0.1387$ and $0.0872$ respectively), with
PC4, PC3, PC5, and PC6 each contributing modestly.  PCs 9 and beyond are nearly
irrelevant for the MLP's predictions.  This is consistent with the PCA scatter
plot (Figure~\ref{fig:pca_labels}), where most visual separation is already
visible along the first two axes.

\begin{table}[H]
  \centering
  \caption{Top 10 most important PCA dimensions for the MLP (permutation
           importance, averaged over 3 shuffles per PC).}
  \label{tab:perm}
  \begin{tabular}{clr}
    \toprule
    Rank & PC & $\Delta$ Macro-F1 \\
    \midrule
    1  & PC2  & 0.1387 \\
    2  & PC1  & 0.0872 \\
    3  & PC4  & 0.0776 \\
    4  & PC3  & 0.0405 \\
    5  & PC5  & 0.0383 \\
    6  & PC6  & 0.0374 \\
    7  & PC9  & 0.0183 \\
    8  & PC7  & 0.0154 \\
    9  & PC11 & 0.0040 \\
    10 & PC8  & 0.0020 \\
    \bottomrule
  \end{tabular}
\end{table}

The sharp drop-off after PC9 ($\Delta F1 < 0.01$) indicates that although
we retain 50 components, the MLP effectively operates on fewer than 10
dimensions.  This aligns with the scree plot (Figure~\ref{fig:pca_variance})
where individual component variance becomes very small after the first dozen PCs.

% ─────────────────────────────────────────────────────────────────────────────
\section{Analysis and Discussion}
\label{sec:discussion}

\subsection{Why Are Accuracies So High?}

The near-perfect classification results ($>$99\% macro-F1) are consistent with
the biology and prior literature:

\begin{itemize}[leftmargin=2em,itemsep=2pt]
  \item \textbf{Tissue-of-origin dominates expression.}  The six cancer types
        originate from entirely different tissues.  Cancer largely preserves
        tissue-specific transcriptional programmes~\cite{weinstein2013}, so the
        discriminative signal is vastly larger than within-type heterogeneity.
  \item \textbf{Batch correction.}  EBPlusPlus normalisation removes
        sequencing artefacts; remaining variance is predominantly biological.
  \item \textbf{Literature precedent.}  Published benchmarks on TCGA Pan-Cancer
        classification report 90--99\% accuracy using RNA-seq
        features~\cite{weinstein2013,hutter2018}.
  \item \textbf{Feature importance validation.}  The top discriminative genes
        (NKX2-1, PCA3, GATA3, SLC22A2, etc.) are established clinical markers,
        confirming the model is learning biology (Section~\ref{sec:importance}).
\end{itemize}

\subsection{MLP vs.\ Logistic Regression}

The MLP achieves a 0.54\,pp test-F1 advantage over LR and also shows lower
cross-validation variance ($\pm 0.0010$ vs $\pm 0.0027$).  The PCA scatter
plot (Figure~\ref{fig:pca_labels}) already shows clean two-dimensional
separation; the MLP's non-linearity provides a minor benefit only at the
finer class boundaries (e.g., BRCA/PRAD overlap in PC1).  Both CV and
permutation importance confirm that the first few PCs carry the bulk of the
discriminative signal, leaving limited room for non-linear models to exploit.

\subsection{Unsupervised vs Supervised Performance}

The ARI gap between K-Means (0.794) and supervised classifiers ($>0.99$) is
expected: K-Means has no label information and must discover structure from
geometry alone.  As shown in Figure~\ref{fig:purity}, five of the six clusters
map cleanly to a single cancer type; residual confusion originates almost
entirely from GBM's small and geometrically overlapping cluster.

\subsection{Data Leakage Prevention}

All \texttt{StandardScaler} and PCA fits are performed inside each CV fold's
training partition only.  The test set is never seen until final evaluation,
and the final MLP is trained using a dev-set monitor split (10\% of dev) rather
than the test set for early stopping diagnostics.
This ensures that the reported metrics faithfully estimate generalisation
performance --- a critical design choice for high-dimensional genomic pipelines
where leakage can produce misleadingly optimistic results.

\subsection{Limitations}

\begin{itemize}[leftmargin=2em,itemsep=1pt]
  \item \textbf{GBM sample size.}  GBM's lower LR-F1 (0.96 vs 1.00 for all
        other classes) traces directly to its 168-sample cohort ($7\times$
        fewer than BRCA); future work could address this with data augmentation
        or over-sampling strategies.
  \item \textbf{Task difficulty.}  Classifying across tissue-of-origin is a
        relatively easy task.  Harder problems (within-BRCA PAM50 subtypes,
        survival prediction) would provide a more discriminating benchmark.
  \item \textbf{Architecture search.}  The MLP architecture [512, 256] was
        fixed; a systematic search could improve performance on harder tasks.
  \item \textbf{Interpretability path.}  The current LR gene projection
        ($\tilde{\mathbf{w}}_c = \mathbf{V}_k^\top \mathbf{w}_c$) identifies
        top genes per class but does not account for gene--gene correlations.
        Pathway-level enrichment analysis on the top gene lists (e.g.,
        Gene Ontology or KEGG) would provide a richer biological interpretation
        and is a natural next step.
\end{itemize}

% ─────────────────────────────────────────────────────────────────────────────
\section{Conclusions}
\label{sec:conclusions}

We implemented and evaluated a complete machine learning pipeline for cancer
type discovery and classification from TCGA Pan-Cancer RNA-seq data, and
extended it with biological validation in response to TA feedback:

\begin{enumerate}[leftmargin=2em,itemsep=2pt]
  \item \textbf{PCA + K-Means} (unsupervised) recovers the six cancer types
        with ARI\,=\,0.794 and NMI\,=\,0.849, demonstrating that bulk RNA-seq
        contains strong tissue-of-origin structure independent of any
        supervised signal.
  \item \textbf{Logistic Regression} ($\ell_2$, OvR) achieves test
        macro-F1\,=\,0.9927 (5-fold CV: $0.9968 \pm 0.0027$), confirming near
        linear separability in 50-dimensional PCA space.
  \item \textbf{MLP} [512, 256] achieves test macro-F1\,=\,0.9981 (5-fold CV:
        $0.9983 \pm 0.0010$), a 0.54\,pp gain over LR with noticeably lower
        cross-fold variance.  The CV estimate provides a conservative
        generalisation bound consistent with the test-set result.
  \item \textbf{Feature importance analysis} (LR weight projection, PC
        loadings, MLP permutation importance) confirms that the models learn
        established tissue-specific gene expression programmes rather than
        technical artefacts: top discriminative genes include NKX2-1/TTF-1
        (LUAD), PCA3 (PRAD), GATA3 and Mammaglobin (BRCA), and SLC22A2/OCT2
        (KIRC) --- all clinically validated cancer biomarkers.
\end{enumerate}

% ─────────────────────────────────────────────────────────────────────────────
\section*{Implementation Notes}

All code is available in the project repository.  Running
\texttt{python src/main.py} from the project root regenerates all figures and
analysis into a timestamped directory \texttt{results/runs/run\_YYYYMMDD\_HHMMSS/}.

\begin{itemize}[leftmargin=2em,itemsep=0pt]
  \item \texttt{src/data\_loader.py} --- GDC API download, log transform, variance filter
  \item \texttt{src/pca.py} --- PCA via truncated SVD (NumPy)
  \item \texttt{src/kmeans.py} --- Lloyd's K-Means with multiple restarts (NumPy)
  \item \texttt{src/logistic\_regression.py} --- OvR LR with $\ell_2$ (NumPy)
  \item \texttt{src/mlp.py} --- feedforward MLP (PyTorch)
  \item \texttt{src/evaluate.py} --- metrics, plotting, and feature importance visualisation
  \item \texttt{src/main.py} --- end-to-end pipeline; generates all figures
\end{itemize}

% ─────────────────────────────────────────────────────────────────────────────
\bibliographystyle{plainnat}
\begin{thebibliography}{9}

\bibitem{weinstein2013}
Weinstein, J.~N., et al.\ (2013).
The Cancer Genome Atlas Pan-Cancer analysis project.
\textit{Nature Genetics}, 45(10), 1113--1120.

\bibitem{hutter2018}
Hutter, C., \& Zenklusen, J.~C.\ (2018).
The Cancer Genome Atlas: creating lasting value beyond its data.
\textit{Cell}, 173(2), 283--285.

\bibitem{cimino2013}
Cimino-Mathews, A., et al.\ (2013).
GATA3 expression in breast carcinoma: utility in triple-negative, sarcomatoid,
and metastatic carcinomas.
\textit{Human Pathology}, 44(7), 1341--1349.

\end{thebibliography}

\end{document}
